\documentclass[12pt]{article}
\usepackage{amsmath}
\usepackage{amsfonts}
\usepackage{float}
\usepackage{fancyhdr}
\usepackage{graphicx}
\usepackage{setspace}
\usepackage[colorlinks=true,linkcolor=blue, citecolor=red]{hyperref}
\usepackage{url}
\usepackage[top=.75in, left=.75in, right=.75in, bottom=1in]{geometry}

% For algorithm
\usepackage{algorithm}
\usepackage{algpseudocode}

% ============ CODE ============
\usepackage{listings}
\usepackage{xcolor}
\definecolor{codegreen}{rgb}{0,0.6,0}
\definecolor{codegray}{rgb}{0.5,0.5,0.5}
\definecolor{codepurple}{rgb}{0.58,0,0.82}
\definecolor{backcolour}{rgb}{0.95,0.95,0.92}

% Styling for the code.
\lstdefinestyle{mystyle}{
    backgroundcolor=\color{backcolour},   
    commentstyle=\color{codegreen},
    keywordstyle=\color{magenta},
    numberstyle=\tiny\color{codegray},
    stringstyle=\color{codepurple},
    basicstyle=\ttfamily\footnotesize,
    breakatwhitespace=false,         
    breaklines=true,                 
    captionpos=b,                    
    keepspaces=true,                 
    numbers=left,                    
    numbersep=5pt,                  
    showspaces=false,                
    showstringspaces=false,
    showtabs=false,                  
    tabsize=2
}
\lstset{style=mystyle}

% Table packages and Unicode / Font support
\usepackage{array}
\usepackage{longtable}
\usepackage{enumitem}
\usepackage{booktabs}
\usepackage[utf8]{inputenc}
\DeclareUnicodeCharacter{2264}{\ensuremath{\le}}
\DeclareUnicodeCharacter{2194}{\ensuremath{\leftrightarrow}}
\usepackage[T1]{fontenc}
\usepackage{lmodern}
\graphicspath{{../../../../templates/}{../../prototype/}{figures/}{./}}
\setlength{\emergencystretch}{3em}
\setlist[enumerate]{label=\arabic*.,leftmargin=*}
\setlist[itemize]{leftmargin=*}

% ============ HEADER AND FOOTER ============
\setlength{\headheight}{29.43912pt}

\pagestyle{fancy}
\fancyfoot{}
\fancyfoot[R]{Page \thepage}

\setlength{\parindent}{0pt}
\renewcommand{\baselinestretch}{1.5}

\newcommand{\coursename}{Introduction of Software Engineering}
\newcommand{\reportname}{Studitopia}
\newcommand{\reporttitle}{\href{https://github.com/RoundFloor/Studitopia.git}{Studitopia}}

\newcommand{\studentname}{%
Nguyen Thanh An (24125024)\\
Nguyen Chi Bao (24125025)\\
Tran Nguyen Khanh Binh (24125026)\\
Tran Canh Anh Tuan (24125048)\\
Nguyen Le Quoc Huy (24125100)
}
\newcommand{\teachername}{Ho Tuan Thanh, Nguyen Thi Minh Tuyen}

% Header
\lhead{\reporttitle}
\rhead{
University of Science - VNUHCM\\
\coursename
}

% Footer
\newcommand{\leftfooter}{}
\lfoot{\leftfooter}

% ============ DOCUMENT ============
\begin{document}

\pagenumbering{roman}
\begin{titlepage}
\newcommand{\HRule}{\rule{\linewidth}{0.5mm}}
\centering

\textsc{\LARGE VNU-HCM}\\[0.5cm]
\textsc{\Large University of Science}\\[0.5cm]
\textsc{\large Faculty of Information Technology}\\[0.5cm]
\textsc{Computer Science}\\[0.5cm]

\HRule \\[0.4cm]
{ 
\huge{\bfseries{\reporttitle}}\\[0.5cm]
\large{\bfseries{Posts}}
}\\[0.4cm]
\HRule \\[0.5cm]

\textbf{\large Subject: \coursename}\\[0.5cm]

\begin{minipage}[t]{0.55\textwidth}
\begin{flushleft} \large
\emph{Students:}\\
\studentname
\end{flushleft}
\end{minipage}
~
\begin{minipage}[t]{0.4\textwidth}
\begin{flushright} \large
\emph{Supervising Teachers:} \\
\teachername
\end{flushright}
\end{minipage}\\[1cm]

{\large \today}\\[1cm]

\includegraphics[scale=.20]{img/hcmus-logo.png}\\[1cm] 

\vfill
\end{titlepage}	

\tableofcontents
\pagebreak

\pagenumbering{arabic}
\setcounter{page}{1}

\hypertarget{sec:posts}{}
\addcontentsline{toc}{section}{Posts}
\section*{Posts}
\par\medskip

\textbf{Version:} \texttt{1.1}
\par\medskip

\textbf{Performed by:} [Nguyen Le Quoc Huy] | \textbf{Edited by:} [Nguyen Le Quoc Huy]
\par\medskip

\texttt{Back to functional requirements index}
\par\medskip

\begin{quote}
\textit{Status: \texttt{Complete} = source-backed implementation exists; \texttt{Partial} = supporting code exists but the full requirement is not implemented; \texttt{Not implemented} = no source-backed implementation was found.}
\end{quote}
\par\medskip

\hypertarget{sec:us-14-create-discussion-post}{}
\addcontentsline{toc}{subsubsection}{US-14: Create Discussion Post}
\subsubsection*{US-14: Create Discussion Post}
\par\medskip

\textbf{As a} university student, \textbf{I want to} create a discussion post, \textbf{so that} I can share my academic problem and ask for solutions.
\par\medskip

\begingroup
\small
\renewcommand{\baselinestretch}{1.0}\selectfont
\renewcommand{\arraystretch}{1.15}
\setlength{\tabcolsep}{6pt}
\begin{longtable}{|>{\raggedright\arraybackslash}p{0.21\textwidth}|>{\raggedright\arraybackslash}p{0.49\textwidth}|>{\raggedright\arraybackslash}p{0.12\textwidth}|>{\raggedright\arraybackslash}p{0.10\textwidth}|}
\hline
\textbf{FR} & \textbf{Requirement} & \textbf{Category} & \textbf{Status} \\ \hline
\endhead
\textbf{FR-14.1 (Discussion Post Creation)} & The system shall allow a user to create a discussion post in a selected forum. & UI, Frontend, Backend & Partial \\ \hline
\textbf{FR-14.2 (Discussion Content)} & The discussion post shall support a title, text content, LaTeX, images, and tags. & UI, Frontend & Not implemented \\ \hline
\textbf{FR-14.3 (Post Category)} & The system shall allow the user to mark the post as a question or a general discussion. & UI, Frontend, Backend & Partial \\ \hline
\textbf{FR-14.4 (Post Forum Context)} & The system shall associate the post with the selected university, major, and course context. & Backend, Database & Not implemented \\ \hline
\textbf{FR-14.5 (Post Validation)} & The system shall validate the post content before submission. & Frontend, Backend & Not implemented \\ \hline
\textbf{FR-14.6 (Post Publication)} & The system shall publish the post and display it in the corresponding forum post list. & Backend, Database & Not implemented \\ \hline
\end{longtable}
\endgroup
\par\medskip

\hypertarget{sec:us-15-create-material-post}{}
\addcontentsline{toc}{subsubsection}{US-15: Create Material Post}
\subsubsection*{US-15: Create Material Post}
\par\medskip

\textbf{As a} university student, \textbf{I want to} create a material post, \textbf{so that} I can share learning resources with other students.
\par\medskip

\begingroup
\small
\renewcommand{\baselinestretch}{1.0}\selectfont
\renewcommand{\arraystretch}{1.15}
\setlength{\tabcolsep}{6pt}
\begin{longtable}{|>{\raggedright\arraybackslash}p{0.16\textwidth}|>{\raggedright\arraybackslash}p{0.55\textwidth}|>{\raggedright\arraybackslash}p{0.10\textwidth}|>{\raggedright\arraybackslash}p{0.10\textwidth}|}
\hline
\textbf{FR} & \textbf{Requirement} & \textbf{Category} & \textbf{Status} \\ \hline
\endhead
\textbf{FR-15.1 (Material Post Creation)} & The system shall allow a user to create a material post in a selected forum. & UI, Frontend, Backend & Partial \\ \hline
\textbf{FR-15.2 (Material Content)} & The material post shall support a title, description, files, and subject tags. & UI, Frontend & Not implemented \\ \hline
\textbf{FR-15.3 (File Validation)} & The system shall validate uploaded files for type and size before accepting them. & Frontend, Backend & Partial \\ \hline
\textbf{FR-15.4 (Material Forum Context)} & The system shall associate the material with the selected university, major, and course context. & Backend, Database & Not implemented \\ \hline
\textbf{FR-15.5 (Material Access)} & The system shall make published materials accessible to users who can view the forum. & Backend & Not implemented \\ \hline
\textbf{FR-15.6 (Material Listing)} & The system shall display the material post in the corresponding forum and in the learning resource library when applicable. & UI, Frontend, Backend & Not implemented \\ \hline
\end{longtable}
\endgroup
\par\medskip

\hypertarget{sec:us-16-create-functional-post}{}
\addcontentsline{toc}{subsubsection}{US-16: Create Functional Post}
\subsubsection*{US-16: Create Functional Post}
\par\medskip

\textbf{As a} university student, \textbf{I want to} create a functional post with an interactive academic tool, \textbf{so that} I can help others practice course content.
\par\medskip

\begingroup
\small
\renewcommand{\baselinestretch}{1.0}\selectfont
\renewcommand{\arraystretch}{1.15}
\setlength{\tabcolsep}{6pt}
\begin{longtable}{|>{\raggedright\arraybackslash}p{0.18\textwidth}|>{\raggedright\arraybackslash}p{0.53\textwidth}|>{\raggedright\arraybackslash}p{0.10\textwidth}|>{\raggedright\arraybackslash}p{0.10\textwidth}|}
\hline
\textbf{FR} & \textbf{Requirement} & \textbf{Category} & \textbf{Status} \\ \hline
\endhead
\textbf{FR-16.1 (Functional Post Creation)} & The system shall allow a user to create a functional post in a selected forum. & UI, Frontend, Backend & Partial \\ \hline
\textbf{FR-16.2 (Functional Content)} & The functional post shall support a title and an interactive tool. & UI, Frontend & Not implemented \\ \hline
\textbf{FR-16.3 (Quiz Tool)} & The system shall support an immediate quiz tool with multiple-choice or single-choice questions. & UI, Frontend, Backend & Not implemented \\ \hline
\textbf{FR-16.4 (Prerequisite Tree Tool)} & The system shall support a prerequisite tree tool that can be created manually or generated from a source file. & UI, Frontend, Backend & Not implemented \\ \hline
\textbf{FR-16.5 (Dynamic Computation)} & The system shall support dynamic computation content when configured by the author. & UI, Frontend, Backend & Not implemented \\ \hline
\textbf{FR-16.6 (Functional Interaction)} & The system shall publish the functional post and allow other users to interact with its tool. & Frontend, Backend & Not implemented \\ \hline
\end{longtable}
\endgroup
\par\medskip

\hypertarget{sec:us-17-view-post-detail}{}
\addcontentsline{toc}{subsubsection}{US-17: View Post Detail}
\subsubsection*{US-17: View Post Detail}
\par\medskip

\textbf{As a} university student, \textbf{I want to} view posts and their details, \textbf{so that} I can read shared content and find relevant information.
\par\medskip

\begingroup
\small
\renewcommand{\baselinestretch}{1.0}\selectfont
\renewcommand{\arraystretch}{1.15}
\setlength{\tabcolsep}{6pt}
\begin{longtable}{|>{\raggedright\arraybackslash}p{0.15\textwidth}|>{\raggedright\arraybackslash}p{0.57\textwidth}|>{\raggedright\arraybackslash}p{0.10\textwidth}|>{\raggedright\arraybackslash}p{0.10\textwidth}|}
\hline
\textbf{FR} & \textbf{Requirement} & \textbf{Category} & \textbf{Status} \\ \hline
\endhead
\textbf{FR-17.1 (Full Post Content)} & The system shall display the full content of a selected post. & UI, Frontend, Backend & Partial \\ \hline
\textbf{FR-17.2 (Post Metadata)} & The system shall display the author, post type, forum context, timestamp, vote score, and tags. & UI, Frontend & Complete \\ \hline
\textbf{FR-17.3 (Discussion Display)} & The system shall display comments, answers, and the marked true response for the post. & UI, Frontend, Backend & Not implemented \\ \hline
\textbf{FR-17.4 (Permission-Based Actions)} & The system shall display available actions such as vote, comment, save, report, and mark true response according to the user's permissions. & UI, Frontend, Backend & Not implemented \\ \hline
\textbf{FR-17.5 (Attachments and Tools)} & The system shall display material attachments and functional tools in the post detail when present. & UI, Frontend & Not implemented \\ \hline
\end{longtable}
\endgroup
\par\medskip

\hypertarget{sec:us-18-comment-and-answer}{}
\addcontentsline{toc}{subsubsection}{US-18: Comment and Answer}
\subsubsection*{US-18: Comment and Answer}
\par\medskip

\textbf{As a} university student, \textbf{I want to} comment on posts and answer questions, \textbf{so that} I can discuss ideas and help other students.
\par\medskip

\begingroup
\small
\renewcommand{\baselinestretch}{1.0}\selectfont
\renewcommand{\arraystretch}{1.15}
\setlength{\tabcolsep}{6pt}
\begin{longtable}{|>{\raggedright\arraybackslash}p{0.21\textwidth}|>{\raggedright\arraybackslash}p{0.48\textwidth}|>{\raggedright\arraybackslash}p{0.12\textwidth}|>{\raggedright\arraybackslash}p{0.10\textwidth}|}
\hline
\textbf{FR} & \textbf{Requirement} & \textbf{Category} & \textbf{Status} \\ \hline
\endhead
\textbf{FR-18.1 (Comment Creation)} & The system shall allow a user to comment on a discussion post. & UI, Frontend, Backend & Not implemented \\ \hline
\textbf{FR-18.2 (Answer Creation)} & The system shall allow a user to answer a question post. & UI, Frontend, Backend & Not implemented \\ \hline
\textbf{FR-18.3 (Nested Replies)} & The system shall support nested replies where applicable. & UI, Frontend, Backend & Not implemented \\ \hline
\textbf{FR-18.4 (Own Content Management)} & The system shall allow users to edit or delete their own comments and answers. & UI, Frontend, Backend & Not implemented \\ \hline
\textbf{FR-18.5 (Comment Ordering)} & The system shall display comments and answers in chronological or threaded order. & UI, Frontend & Not implemented \\ \hline
\textbf{FR-18.6 (Author Notification)} & The system shall notify the post author when a new comment or answer is added. & Backend & Not implemented \\ \hline
\end{longtable}
\endgroup
\par\medskip

\hypertarget{sec:us-19-vote-on-posts-and-comments}{}
\addcontentsline{toc}{subsubsection}{US-19: Vote on Posts and Comments}
\subsubsection*{US-19: Vote on Posts and Comments}
\par\medskip

\textbf{As a} university student, \textbf{I want to} vote on posts and comments, \textbf{so that} useful academic content is easier to discover.
\par\medskip

\begingroup
\small
\renewcommand{\baselinestretch}{1.0}\selectfont
\renewcommand{\arraystretch}{1.15}
\setlength{\tabcolsep}{6pt}
\begin{longtable}{|>{\raggedright\arraybackslash}p{0.21\textwidth}|>{\raggedright\arraybackslash}p{0.46\textwidth}|>{\raggedright\arraybackslash}p{0.15\textwidth}|>{\raggedright\arraybackslash}p{0.10\textwidth}|}
\hline
\textbf{FR} & \textbf{Requirement} & \textbf{Category} & \textbf{Status} \\ \hline
\endhead
\textbf{FR-19.1 (Post Voting)} & The system shall allow a user to upvote or downvote a post. & UI, Frontend, Backend & Not implemented \\ \hline
\textbf{FR-19.2 (Comment Voting)} & The system shall allow a user to upvote or downvote a comment or answer. & UI, Frontend, Backend & Not implemented \\ \hline
\textbf{FR-19.3 (Duplicate Vote Prevention)} & The system shall prevent a user from voting more than once on the same post or comment. & Frontend, Backend, Database & Not implemented \\ \hline
\textbf{FR-19.4 (Vote Change)} & The system shall allow a user to change or remove their vote. & UI, Frontend, Backend & Not implemented \\ \hline
\textbf{FR-19.5 (Score Update)} & The system shall update the vote score and ranking after a vote is recorded. & Backend, Database & Not implemented \\ \hline
\end{longtable}
\endgroup
\par\medskip

\hypertarget{sec:us-20-mark-true-response}{}
\addcontentsline{toc}{subsubsection}{US-20: Mark True Response}
\subsubsection*{US-20: Mark True Response}
\par\medskip

\textbf{As a} post author, \textbf{I want to} mark a comment or answer as the true response, \textbf{so that} other students can quickly find the accepted solution.
\par\medskip

\begingroup
\small
\renewcommand{\baselinestretch}{1.0}\selectfont
\renewcommand{\arraystretch}{1.15}
\setlength{\tabcolsep}{6pt}
\begin{longtable}{|>{\raggedright\arraybackslash}p{0.16\textwidth}|>{\raggedright\arraybackslash}p{0.51\textwidth}|>{\raggedright\arraybackslash}p{0.14\textwidth}|>{\raggedright\arraybackslash}p{0.10\textwidth}|}
\hline
\textbf{FR} & \textbf{Requirement} & \textbf{Category} & \textbf{Status} \\ \hline
\endhead
\textbf{FR-20.1 (Mark Response)} & The system shall allow the author of a question post to mark one comment or answer as the true response. & UI, Frontend, Backend & Not implemented \\ \hline
\textbf{FR-20.2 (Solved Status)} & The system shall mark the post as solved when a true response is selected. & Frontend, Backend, Database & Not implemented \\ \hline
\textbf{FR-20.3 (Prominent Display)} & The system shall display the true response prominently under the post. & UI, Frontend & Not implemented \\ \hline
\textbf{FR-20.4 (Change Response)} & The system shall allow the post author to change the true response when necessary. & UI, Frontend, Backend & Not implemented \\ \hline
\textbf{FR-20.5 (Answer Notification)} & The system shall notify the author of the selected answer that their answer was accepted. & Backend & Not implemented \\ \hline
\end{longtable}
\endgroup
\par\medskip

\hypertarget{sec:us-21-save-post}{}
\addcontentsline{toc}{subsubsection}{US-21: Save Post}
\subsubsection*{US-21: Save Post}
\par\medskip

\textbf{As a} university student, \textbf{I want to} save posts, \textbf{so that} I can return to useful discussions and materials later.
\par\medskip

\begingroup
\small
\renewcommand{\baselinestretch}{1.0}\selectfont
\renewcommand{\arraystretch}{1.15}
\setlength{\tabcolsep}{6pt}
\begin{longtable}{|>{\raggedright\arraybackslash}p{0.19\textwidth}|>{\raggedright\arraybackslash}p{0.50\textwidth}|>{\raggedright\arraybackslash}p{0.13\textwidth}|>{\raggedright\arraybackslash}p{0.10\textwidth}|}
\hline
\textbf{FR} & \textbf{Requirement} & \textbf{Category} & \textbf{Status} \\ \hline
\endhead
\textbf{FR-21.1 (Save Post)} & The system shall allow a user to save a post. & UI, Frontend, Backend & Not implemented \\ \hline
\textbf{FR-21.2 (Saved List)} & The system shall provide a list of the user's saved posts. & UI, Frontend, Backend & Not implemented \\ \hline
\textbf{FR-21.3 (Remove Saved Post)} & The system shall allow a user to remove a post from their saved list. & UI, Frontend, Backend & Not implemented \\ \hline
\textbf{FR-21.4 (Saved Persistence)} & The system shall keep saved posts available even when the user changes forum context. & Backend, Database & Not implemented \\ \hline
\end{longtable}
\endgroup
\par\medskip

\hypertarget{sec:us-22-report-and-moderate-forum-content}{}
\addcontentsline{toc}{subsubsection}{US-22: Report and Moderate Forum Content}
\subsubsection*{US-22: Report and Moderate Forum Content}
\par\medskip

\textbf{As a} university student or moderator, \textbf{I want to} report and moderate forum content, \textbf{so that} the forum remains safe and academically useful.
\par\medskip

\begingroup
\small
\renewcommand{\baselinestretch}{1.0}\selectfont
\renewcommand{\arraystretch}{1.15}
\setlength{\tabcolsep}{6pt}
\begin{longtable}{|>{\raggedright\arraybackslash}p{0.17\textwidth}|>{\raggedright\arraybackslash}p{0.54\textwidth}|>{\raggedright\arraybackslash}p{0.10\textwidth}|>{\raggedright\arraybackslash}p{0.10\textwidth}|}
\hline
\textbf{FR} & \textbf{Requirement} & \textbf{Category} & \textbf{Status} \\ \hline
\endhead
\textbf{FR-22.1 (Report Content)} & The system shall allow a user to report a post, comment, or answer. & UI, Frontend, Backend & Not implemented \\ \hline
\textbf{FR-22.2 (Report Categories)} & The system shall provide report categories such as spam, harassment, inappropriate content, and academic dishonesty. & UI, Frontend & Not implemented \\ \hline
\textbf{FR-22.3 (Moderation Queue)} & The system shall route reports to a moderation queue for review. & Backend & Not implemented \\ \hline
\textbf{FR-22.4 (Moderator Actions)} & The system shall allow moderators to hide, delete, restore, or ban content and users according to forum policy. & UI, Frontend, Backend & Not implemented \\ \hline
\textbf{FR-22.5 (Moderation Record)} & The system shall record moderation actions and notify the affected user when appropriate. & Backend, Database & Not implemented \\ \hline
\textbf{FR-22.6 (Restriction Enforcement)} & The system shall prevent moderated users from bypassing restrictions through new posts or comments. & Backend & Not implemented \\ \hline
\end{longtable}
\endgroup
\par\medskip

\par\medskip
\hypertarget{sec:ui-prototype-specimens}{}
\addcontentsline{toc}{subsection}{UI Prototype Specimens}
\subsection*{UI Prototype Specimens}
\par\medskip
\begin{figure}[H]
\centering
\includegraphics[width=0.75\textwidth,height=0.42\textheight,keepaspectratio]{"../../prototype/create_post_desktop.png"}
\par\vspace{2pt}
{\footnotesize \textbf{UI Prototype:} Create post - Desktop}
\end{figure}
\par\smallskip
\begin{figure}[H]
\centering
\includegraphics[width=0.75\textwidth,height=0.42\textheight,keepaspectratio]{"../../prototype/create_post_discussion.png"}
\par\vspace{2pt}
{\footnotesize \textbf{UI Prototype:} Create post - Discussion}
\end{figure}
\par\smallskip
\begin{figure}[H]
\centering
\includegraphics[width=0.75\textwidth,height=0.42\textheight,keepaspectratio]{"../../prototype/create_post_functional_quiz.png"}
\par\vspace{2pt}
{\footnotesize \textbf{UI Prototype:} Create post - Functional - Quiz}
\end{figure}
\par\smallskip
\begin{figure}[H]
\centering
\includegraphics[width=0.75\textwidth,height=0.42\textheight,keepaspectratio]{"../../prototype/create_post_material.png"}
\par\vspace{2pt}
{\footnotesize \textbf{UI Prototype:} Create post - Material}
\end{figure}
\par\smallskip
\begin{figure}[H]
\centering
\includegraphics[width=0.75\textwidth,height=0.42\textheight,keepaspectratio]{"../../prototype/create_post_prerequisite_tree_code.png"}
\par\vspace{2pt}
{\footnotesize \textbf{UI Prototype:} Create post - Prerequisite tree - Code}
\end{figure}
\par\smallskip
\begin{figure}[H]
\centering
\includegraphics[width=0.75\textwidth,height=0.42\textheight,keepaspectratio]{"../../prototype/create_post_prerequisite_tree_drag_and_drop.png"}
\par\vspace{2pt}
{\footnotesize \textbf{UI Prototype:} Create post - Prerequisite tree - Drag and drop}
\end{figure}
\par\smallskip
\begin{figure}[H]
\centering
\includegraphics[width=0.75\textwidth,height=0.42\textheight,keepaspectratio]{"../../prototype/create_post_prerequisite_tree_generate_from_file.png"}
\par\vspace{2pt}
{\footnotesize \textbf{UI Prototype:} Create post - Prerequisite tree - Generate from file}
\end{figure}
\par\smallskip
\begin{figure}[H]
\centering
\includegraphics[width=0.75\textwidth,height=0.42\textheight,keepaspectratio]{"../../prototype/create_post_validation_and_publication_states.png"}
\par\vspace{2pt}
{\footnotesize \textbf{UI Prototype:} Create post - Validation and publication states}
\end{figure}
\par\smallskip
\begin{figure}[H]
\centering
\includegraphics[width=0.75\textwidth,height=0.42\textheight,keepaspectratio]{"../../prototype/dijkstra_question_desktop.png"}
\par\vspace{2pt}
{\footnotesize \textbf{UI Prototype:} Dijkstra question - Desktop}
\end{figure}
\par\smallskip
\begin{figure}[H]
\centering
\includegraphics[width=0.75\textwidth,height=0.42\textheight,keepaspectratio]{"../../prototype/prerequisite_tree_validation_and_recovery_states.png"}
\par\vspace{2pt}
{\footnotesize \textbf{UI Prototype:} Prerequisite tree - Validation and recovery states}
\end{figure}
\par\smallskip
\begin{figure}[H]
\centering
\includegraphics[width=0.75\textwidth,height=0.42\textheight,keepaspectratio]{"../../prototype/published_prerequisite_tree_reader_simulation.png"}
\par\vspace{2pt}
{\footnotesize \textbf{UI Prototype:} Published prerequisite tree - Reader simulation}
\end{figure}
\par\smallskip

\end{document}
